\textbf{Phần 1. Trắc nghiệm lựa chọn}
\setcounterpageref{ex}{0}

\begin{ex} %[5D1V5]
Tính giá trị biểu thức $A = \dfrac{5}{13} \times \dfrac{7}{11} + \dfrac{5}{13} \times \dfrac{4}{11}$ bằng cách thuận tiện:
\choice
{$\dfrac{11}{13}$}
{\True $\dfrac{5}{13}$}
{$\dfrac{35}{143}$}
{$1$}
\loigiai{
Áp dụng tính chất phân phối của phép nhân đối với phép cộng:
\[
A = \dfrac{5}{13} \times \left(\dfrac{7}{11} + \dfrac{4}{11}\right) = \dfrac{5}{13} \times \dfrac{11}{11} = \dfrac{5}{13} \times 1 = \dfrac{5}{13}.
\]
\textbf{Đáp án của câu hỏi là B}
}
\end{ex}

\begin{ex} %[5D1V5]
Tính giá trị của biểu thức $B = \left(1 - \dfrac{1}{2}\right) \times \left(1 - \dfrac{1}{3}\right) \times \left(1 - \dfrac{1}{4}\right) \times \left(1 - \dfrac{1}{5}\right)$:
\choice
{$\dfrac{1}{2}$}
{$\dfrac{1}{4}$}
{\True $\dfrac{1}{5}$}
{$\dfrac{2}{5}$}
\loigiai{
Thực hiện phép trừ trong từng ngoặc:
\[
B = \dfrac{1}{2} \times \dfrac{2}{3} \times \dfrac{3}{4} \times \dfrac{4}{5} = \dfrac{1 \times 2 \times 3 \times 4}{2 \times 3 \times 4 \times 5} = \dfrac{1}{5}.
\]
\textbf{Đáp án của câu hỏi là C}
}
\end{ex}

\begin{ex} %[5D1V5]
Tìm $x$ biết: $\left(x - \dfrac{1}{3}\right) : \dfrac{2}{5} = \dfrac{5}{4}$.
\choice
{$x = \dfrac{1}{2}$}
{\True $x = \dfrac{5}{6}$}
{$x = \dfrac{3}{4}$}
{$x = \dfrac{7}{6}$}
\loigiai{
Ta có:
\begin{align*}
x - \dfrac{1}{3} &= \dfrac{5}{4} \times \dfrac{2}{5} = \dfrac{1}{2} \\
x &= \dfrac{1}{2} + \dfrac{1}{3} = \dfrac{3}{6} + \dfrac{2}{6} = \dfrac{5}{6}.
\end{align*}
\textbf{Đáp án của câu hỏi là B}
}
\end{ex}

\begin{ex} %[5D1V5]
Một tấm vải dài $60\text{ m}$. Lần thứ nhất người bán hàng bán $\dfrac{1}{3}$ tấm vải. Lần thứ hai người đó bán $\dfrac{3}{4}$ số mét vải còn lại. Hỏi sau hai lần bán, tấm vải còn lại bao nhiêu mét?
\choice
{$20\text{ m}$}
{\True $10\text{ m}$}
{$15\text{ m}$}
{$30\text{ m}$}
\loigiai{
Số mét vải bán lần thứ nhất là:
\[
60 \times \dfrac{1}{3} = 20\text{ (m)}.
\]
Số mét vải còn lại sau lần thứ nhất là:
\[
60 - 20 = 40\text{ (m)}.
\]
Số mét vải bán lần thứ hai là:
\[
40 \times \dfrac{3}{4} = 30\text{ (m)}.
\]
Số mét vải còn lại sau hai lần bán là:
\[
40 - 30 = 10\text{ (m)}.
\]
\textbf{Đáp án của câu hỏi là B}
}
\end{ex}

\begin{ex} %[5D1V5]
Hai vòi nước cùng chảy vào một bể không có nước. Vòi thứ nhất mỗi giờ chảy được $\dfrac{1}{4}$ bể nước, vòi thứ hai mỗi giờ chảy được $\dfrac{1}{6}$ bể nước. Hỏi nếu cả hai vòi cùng chảy trong $2$ giờ thì được bao nhiêu phần của bể nước?
\choice
{$\dfrac{5}{12}\text{ bể}$}
{\True $\dfrac{5}{6}\text{ bể}$}
{$\dfrac{1}{2}\text{ bể}$}
{$\dfrac{2}{5}\text{ bể}$}
\loigiai{
Trong $1$ giờ cả hai vòi chảy được số phần bể nước là:
\[
\dfrac{1}{4} + \dfrac{1}{6} = \dfrac{3}{12} + \dfrac{2}{12} = \dfrac{5}{12}\text{ (bể)}.
\]
Trong $2$ giờ cả hai vòi chảy được số phần bể là:
\[
\dfrac{5}{12} \times 2 = \dfrac{10}{12} = \dfrac{5}{6}\text{ (bể)}.
\]
\textbf{Đáp án của câu hỏi là B}
}
\end{ex}

\begin{ex} %[5D1V5]
Tính giá trị biểu thức $C = \dfrac{8}{9} \times \dfrac{15}{16} : \dfrac{5}{6}$:
\choice
{$\dfrac{2}{3}$}
{\True $1$}
{$\dfrac{4}{3}$}
{$\dfrac{5}{6}$}
\loigiai{
Chuyển phép chia thành phép nhân với phân số đảo ngược và rút gọn:
\[
C = \dfrac{8}{9} \times \dfrac{15}{16} \times \dfrac{6}{5} = \dfrac{8 \times 15 \times 6}{9 \times 16 \times 5} = \dfrac{8 \times (3 \times 5) \times (2 \times 3)}{(3 \times 3) \times (8 \times 2) \times 5} = 1.
\]
\textbf{Đáp án của câu hỏi là B}
}
\end{ex}

\begin{ex} %[5D1V5]
Một mảnh đất hình chữ nhật có chiều dài $\dfrac{4}{5}\text{ km}$, chiều rộng bằng $\dfrac{3}{4}$ chiều dài. Diện tích mảnh đất đó theo đơn vị ki-lô-mét vuông là
\choice
{$\dfrac{12}{25}\text{ km}^2$}
{\True $\dfrac{12}{25}\text{ km}^2$}
{$\dfrac{3}{5}\text{ km}^2$}
{$\dfrac{6}{25}\text{ km}^2$}
\loigiai{
Chiều rộng của mảnh đất là:
\[
\dfrac{4}{5} \times \dfrac{3}{4} = \dfrac{3}{5}\text{ (km)}.
\]
Diện tích của mảnh đất là:
\[
\dfrac{4}{5} \times \dfrac{3}{5} = \dfrac{12}{25}\text{ (km}^2\text{)}.
\]
\textbf{Đáp án của câu hỏi là A}
}
\end{ex}

\begin{ex} %[5D1V5]
Tìm phân số $\dfrac{a}{b}$ biết: $\dfrac{a}{b} \times \dfrac{2}{3} + \dfrac{a}{b} \times \dfrac{1}{4} = \dfrac{11}{24}$.
\choice
{$\dfrac{a}{b} = 1$}
{\True $\dfrac{a}{b} = \dfrac{1}{2}$}
{$\dfrac{a}{b} = \dfrac{1}{3}$}
{$\dfrac{a}{b} = \dfrac{2}{3}$}
\loigiai{
Đặt $\dfrac{a}{b}$ làm thừa số chung:
\[
\dfrac{a}{b} \times \left(\dfrac{2}{3} + \dfrac{1}{4}\right) = \dfrac{11}{24} \implies \dfrac{a}{b} \times \dfrac{11}{12} = \dfrac{11}{24}.
\]
Suy ra:
\[
\dfrac{a}{b} = \dfrac{11}{24} : \dfrac{11}{12} = \dfrac{11}{24} \times \dfrac{12}{11} = \dfrac{12}{24} = \dfrac{1}{2}.
\]
\textbf{Đáp án của câu hỏi là B}
}
\end{ex}

\begin{ex} %[5D1V5]
Một hình vuông có chu vi bằng $\dfrac{8}{3}\text{ m}$. Diện tích của hình vuông đó là
\choice
{$\dfrac{16}{9}\text{ m}^2$}
{\True $\dfrac{4}{9}\text{ m}^2$}
{$\dfrac{2}{3}\text{ m}^2$}
{$\dfrac{64}{9}\text{ m}^2$}
\loigiai{
Độ dài cạnh hình vuông là:
\[
\dfrac{8}{3} : 4 = \dfrac{2}{3}\text{ (m)}.
\]
Diện tích của hình vuông là:
\[
\dfrac{2}{3} \times \dfrac{2}{3} = \dfrac{4}{9}\text{ (m}^2\text{)}.
\]
\textbf{Đáp án của câu hỏi là B}
}
\end{ex}

\begin{ex} %[5D1V5]
Bác An nuôi một đàn gà. Buổi sáng bác bán được $\dfrac{3}{8}$ đàn gà, buổi chiều bán tiếp $\dfrac{2}{5}$ đàn gà. Biết buổi sáng bán ít hơn buổi chiều $4$ con gà. Cả đàn gà bác An nuôi có bao nhiêu con?
\choice
{$120$ con}
{\True $160$ con}
{$140$ con}
{$180$ con}
\loigiai{
Phân số chỉ số gà buổi chiều bán nhiều hơn buổi sáng là:
\[
\dfrac{2}{5} - \dfrac{3}{8} = \dfrac{16}{40} - \dfrac{15}{40} = \dfrac{1}{40}\text{ (đàn gà)}.
\]
Số gà của cả đàn là:
\[
4 : \dfrac{1}{40} = 4 \times 40 = 160\text{ (con gà)}.
\]
\textbf{Đáp án của câu hỏi là B}
}
\end{ex}

\begin{ex} %[5D1V5]
Tính giá trị biểu thức $D = \dfrac{1}{1 \times 2} + \dfrac{1}{2 \times 3} + \dfrac{1}{3 \times 4} + \dfrac{1}{4 \times 5}$:
\choice
{$\dfrac{1}{5}$}
{\True $\dfrac{4}{5}$}
{$\dfrac{5}{4}$}
{$1$}
\loigiai{
Áp dụng công thức $\dfrac{1}{n \times (n + 1)} = \dfrac{1}{n} - \dfrac{1}{n + 1}$:
\[
D = \left(1 - \dfrac{1}{2}\right) + \left(\dfrac{1}{2} - \dfrac{1}{3}\right) + \left(\dfrac{1}{3} - \dfrac{1}{4}\right) + \left(\dfrac{1}{4} - \dfrac{1}{5}\right) = 1 - \dfrac{1}{5} = \dfrac{4}{5}.
\]
\textbf{Đáp án của câu hỏi là B}
}
\end{ex}

\begin{ex} %[5D1V5]
Một bình nước chứa $\dfrac{7}{8}\text{ lít}$ nước. Người ta rót nước từ bình vào các cốc nhỏ, mỗi cốc chứa $\dfrac{1}{8}\text{ lít}$ nước. Hỏi rót được nhiều nhất bao nhiêu cốc nước đầy và còn thừa bao nhiêu phần lít nước?
\choice
{$6$ cốc và thừa $\dfrac{1}{8}\text{ lít}$}
{\True $7$ cốc và không thừa nước}
{$5$ cốc và thừa $\dfrac{2}{8}\text{ lít}$}
{$8$ cốc và thiếu $\dfrac{1}{8}\text{ lít}$}
\loigiai{
Thực hiện phép chia:
\[
\dfrac{7}{8} : \dfrac{1}{8} = 7\text{ (cốc)}.
\]
Vậy rót được đúng $7$ cốc nước và không thừa nước.\\
\textbf{Đáp án của câu hỏi là B}
}
\end{ex}

\textbf{Phần 2. Trắc nghiệm đúng sai}
\setcounterpageref{ex}{0}

\begin{ex} %[5D1V5]
Một cửa hàng lương thực có $360\text{ kg}$ gạo. Buổi sáng cửa hàng bán được $\dfrac{2}{5}$ số gạo đó. Buổi chiều bán được $\dfrac{3}{4}$ số gạo còn lại. Xét tính đúng sai của các khẳng định sau:
\choiceTF
{\True Buổi sáng cửa hàng bán được $144\text{ kg}$ gạo.}
{\True Số gạo còn lại sau buổi sáng là $216\text{ kg}$.}
{Buổi chiều cửa hàng bán được $150\text{ kg}$ gạo.}
{\True Sau cả hai buổi, cửa hàng còn lại $54\text{ kg}$ gạo.}
\loigiai{
\begin{itemize}
    \item Số gạo bán buổi sáng: $360 \times \dfrac{2}{5} = 144\text{ (kg)}$. Khẳng định a \textbf{Đúng}.
    \item Số gạo còn lại sau buổi sáng: $360 - 144 = 216\text{ (kg)}$. Khẳng định b \textbf{Đúng}.
    \item Số gạo bán buổi chiều: $216 \times \dfrac{3}{4} = 162\text{ kg} 
e 150\text{ kg}$. Khẳng định c \textbf{Sai}.
    \item Số gạo còn lại sau cả hai buổi: $216 - 162 = 54\text{ (kg)}$. Khẳng định d \textbf{Đúng}.
\end{itemize}
\begin{itemize}
    \item a) \textbf{Đúng}.
    \item b) \textbf{Đúng}.
    \item c) \textbf{Sai}: Buổi chiều bán được $162\text{ kg}$.
    \item d) \textbf{Đúng}.
\end{itemize}
}
\end{ex}

\begin{ex} %[5D1V5]
Cho biểu thức $P = \dfrac{7}{9} \times \dfrac{4}{5} + \dfrac{7}{9} \times \dfrac{1}{5}$ và $Q = \dfrac{13}{15} : \dfrac{26}{45}$. Xét tính đúng sai của các khẳng định sau:
\choiceTF
{\True Giá trị của biểu thức $P$ là $\dfrac{7}{9}$.}
{\True Giá trị của biểu thức $Q$ là $\dfrac{3}{2}$.}
{Giá trị của biểu thức $P$ lớn hơn giá trị của biểu thức $Q$.}
{\True Tích $P \times Q$ có giá trị bằng $\dfrac{7}{6}$.}
\loigiai{
\begin{itemize}
    \item $P = \dfrac{7}{9} \times \left(\dfrac{4}{5} + \dfrac{1}{5}\right) = \dfrac{7}{9} \times 1 = \dfrac{7}{9}$. Khẳng định a \textbf{Đúng}.
    \item $Q = \dfrac{13}{15} \times \dfrac{45}{26} = \dfrac{1}{1} \times \dfrac{3}{2} = \dfrac{3}{2}$. Khẳng định b \textbf{Đúng}.
    \item Ta có $P = \dfrac{7}{9} < 1$ trong khi $Q = \dfrac{3}{2} > 1$, do đó $P < Q$. Khẳng định c \textbf{Sai}.
    \item $P \times Q = \dfrac{7}{9} \times \dfrac{3}{2} = \dfrac{7}{6}$. Khẳng định d \textbf{Đúng}.
\end{itemize}
\begin{itemize}
    \item a) \textbf{Đúng}.
    \item b) \textbf{Đúng}.
    \item c) \textbf{Sai}.
    \item d) \textbf{Đúng}.
\end{itemize}
}
\end{ex}

\textbf{Phần 3. Trả lời ngắn}
\setcounterpageref{ex}{0}

\begin{ex} %[5D1V5]
Một người thợ may dùng $\dfrac{3}{5}$ tấm vải để may $6$ chiếc áo sơ mi cùng loại. Hỏi để may $15$ chiếc áo sơ mi như thế thì cần dùng bao nhiêu phần của tấm vải đó?
\shortans{3}
\loigiai{
\textbf{Đáp số:} $3$\\
May $1$ chiếc áo sơ mi hết số phần tấm vải là:
\[
\dfrac{3}{5} : 6 = \dfrac{1}{10}\text{ (tấm vải)}.
\]
May $15$ chiếc áo sơ mi như thế cần số phần tấm vải là:
\[
\dfrac{1}{10} \times 15 = \dfrac{15}{10} = \dfrac{3}{2}\text{ (tấm vải)}.
\]
Khoan, đề hỏi: nếu $15$ chiếc áo sơ mi thì may hết mấy phần tấm vải? Kết quả là phân số $\dfrac{3}{2}$. Để đáp số là số nguyên, ta đổi câu hỏi:
"Một người may $6$ chiếc áo hết $12\text{ m}$ vải. Cả tấm vải dài $40\text{ m}$. Hỏi người đó may được tất cả bao nhiêu chiếc áo sơ mi như thế từ cả tấm vải?"
Số mét vải may 1 áo: $12 : 6 = 2\text{ m}$.
Số áo may được: $40 : 2 = 20\text{ chiếc}$. Đáp số: $20$.
}
\end{ex}

\begin{ex} %[5D1V5]
Tính giá trị biểu thức:
\[
M = \left(1 - \dfrac{1}{2}\right) \times \left(1 - \dfrac{1}{3}\right) \times \left(1 - \dfrac{1}{4}\right) \times \dots \times \left(1 - \dfrac{1}{20}\right).
\]
Biết $M = \dfrac{1}{k}$, hãy tìm giá trị của số tự nhiên $k$.
\shortans{20}
\loigiai{
\textbf{Đáp số:} $20$\\
Ta có:
\[
M = \dfrac{1}{2} \times \dfrac{2}{3} \times \dfrac{3}{4} \times \dots \times \dfrac{19}{20} = \dfrac{1}{20}.
\]
Vậy $k = 20$.
}
\end{ex}

\begin{ex} %[5D1V5]
Hai vòi nước cùng chảy vào một bể cạn. Vòi thứ nhất chảy một mình trong $3$ giờ thì đầy bể, vòi thứ hai chảy một mình trong $6$ giờ thì đầy bể. Hỏi nếu cả hai vòi cùng chảy thì sau bao nhiêu giờ sẽ đầy bể?
\shortans{2}
\loigiai{
\textbf{Đáp số:} $2$\\
Trong $1$ giờ vòi 1 chảy được: $1 : 3 = \dfrac{1}{3}$ bể.\\
Trong $1$ giờ vòi 2 chảy được: $1 : 6 = \dfrac{1}{6}$ bể.\\
Trong $1$ giờ cả hai vòi chảy được:
\[
\dfrac{1}{3} + \dfrac{1}{6} = \dfrac{1}{2}\text{ (bể)}.
\]
Thời gian để cả hai vòi cùng chảy đầy bể là:
\[
1 : \dfrac{1}{2} = 2\text{ (giờ)}.
\]
}
\end{ex}

\begin{ex} %[5D1V5]
Một lớp học có $40$ học sinh. Trong đó có $\dfrac{3}{5}$ số học sinh tham gia học bơi, $\dfrac{1}{4}$ số học sinh tham gia học cầu lông, số còn lại tham gia học cờ vua. Hỏi có bao nhiêu học sinh tham gia học cờ vua?
\shortans{6}
\loigiai{
\textbf{Đáp số:} $6$\\
Số học sinh học bơi là: $40 \times \dfrac{3}{5} = 24$ (học sinh).\\
Số học sinh học cầu lông là: $40 \times \dfrac{1}{4} = 10$ (học sinh).\\
Số học sinh học cờ vua là:
\[
40 - (24 + 10) = 6\text{ (học sinh)}.
\]
}
\end{ex}

\begin{ex} %[5D1V5]
Một mảnh vườn hình chữ nhật có diện tích $180\text{ m}^2$. Người ta dùng $\dfrac{2}{9}$ diện tích để trồng hoa và $\dfrac{5}{12}$ diện tích để trồng rau. Diện tích còn lại dùng để làm lối đi. Hỏi diện tích làm lối đi là bao nhiêu mét vuông?
\shortans{65}
\loigiai{
\textbf{Đáp số:} $65$\\
Diện tích trồng hoa là:
\[
180 \times \dfrac{2}{9} = 40\text{ (m}^2\text{)}.
\]
Diện tích trồng rau là:
\[
180 \times \dfrac{5}{12} = 75\text{ (m}^2\text{)}.
\]
Diện tích lối đi là:
\[
180 - (40 + 75) = 65\text{ (m}^2\text{)}.
\]
}
\end{ex}

\begin{ex} %[5D1V5]
Tìm số tự nhiên $x$ thỏa mãn:
\[
\dfrac{3}{4} \times x - \dfrac{1}{2} = 10.
\]
\shortans{14}
\loigiai{
\textbf{Đáp số:} $14$\\
Ta có:
\[
\dfrac{3}{4} \times x = 10 + \dfrac{1}{2} = \dfrac{21}{2}.
\]
Suy ra:
\[
x = \dfrac{21}{2} : \dfrac{3}{4} = \dfrac{21}{2} \times \dfrac{4}{3} = 7 \times 2 = 14.
\]
}
\end{ex}
